\section{Theory}
\label{sec:theory}

One reading organises everything that follows.  The design is a horizon-one
decision process with a context-dependent action set: at a sampled checkpoint
the controller selects one intervention option from the feasible set, the frozen
continuation $\rho$ then runs to the horizon, and validated completion is the
terminal reward.  Every quantity the branch design estimates is therefore a
Bellman object of the baseline controller that never intervenes.  Writing
$V^\rho$ and $Q^\rho$ for its state- and action-value functions at the sampled
checkpoint, evaluated at the augmented state $X$ of \cref{sec:b-problem}, the
conditional regime value $\mu_g$ is an aggregate of $Q^\rho(\cdot,g)$, baseline
failure \emph{risk} is $1-V^\rho$, and verification advantage and actionable
value are advantages $A^\rho=Q^\rho-V^\rho$; a branch is a Monte-Carlo backup of
$Q^\rho$ from a restored state, and the availability mask is membership in an
option's initiation set.  \Cref{tab:bellman-dictionary} in \cref{app:notation}
gives the full correspondence.  The reading is used for two purposes only.  It
states the retained policy-ranking result in familiar terms --- \cref{thm:ordinal} is the exact
condition under which the failure-risk function $1-V^\rho$ ranks the advantage
function $A^\rho(\cdot,g_v)$; ranking by $V^\rho$ itself is the same comparison
with the sign reversed, and it is failure risk that a monitor reports --- and,
in \cref{sec:bellman}, it makes the one-step restriction a quantified choice
rather than a scoping sentence.  It is not a method: the learner of
\cref{sec:rl} remains a one-step full-information procedure that optimises a
masked policy score rather than regressing an action value, and no sequential
value iteration is proposed.  Nor is it an identification argument.  A Bellman
recursion is an accounting identity for whichever kernel generated the data; it
is \cref{sec:b-data} and the runtime qualification gate, not the recursion, that decide
whether that kernel is the one the branches executed.

Four results.  Each determines a component of the design or a prediction in
\cref{sec:policy,sec:b-experiments}; none requires data; all concern quantities the
branch design already estimates.  The first supplies the sharp null of the
population-level ranking comparison: exactly when a baseline failure-risk score suffices to
allocate a verification budget --- in the reading above, exactly when $1-V^\rho$ ranks
$A^\rho(\cdot,g_v)$ (\cref{thm:ordinal}).  The remaining results are the machinery that makes
acting on that answer legitimate --- a bound on the loss from registering a
representation (\cref{thm:refinement}), a pessimistic gated rule that cannot
make a running system worse (\cref{thm:safe,cor:proposal}), and the simultaneous
certification that prices the whole procedure (\cref{prop:certification}).

The results are individually elementary, and we state that plainly: the
contribution is the object they apply to, together with the observation that
\emph{the one-step restriction makes their hypotheses verifiable rather than
assumed}.  A sequential safe-improvement guarantee of the kind in
\cref{thm:safe} would require uniform confidence validity over the states
reachable under the new policy, which is unknown and demands a positivity
assumption over a sequential action space.  Executing every feasible regime at a
single randomly selected checkpoint removes that problem for the
\emph{action}: conditional on a checkpoint being sampled, every regime feasible
there is executed, so action-assignment positivity holds by construction rather
than by assumption.  It does not promise that every cell of a partition is
sampled: a cell the registered sampling law gives zero or negligible mass
receives no branches, and the gate treats it as uncovered and uncertifiable
rather than as certified by default.

Two things must be kept apart.  Exhaustive execution supplies
\emph{coverage}: every feasible regime is executed at every eligible checkpoint,
so for the execution indicator $E_{ig}$ $\Prb(E_{ig}=1\mid A_{ig}=1)=1$.  It does
not supply \emph{statistical} validity: the uniform confidence bounds
\cref{thm:safe} requires are the gate family of \cref{eq:gate-family}, which
depends on a registered finite partition of the representation space, the sample
size, and a simultaneous union bound over the gate and certification families on
held-out calibration data.  \textbf{Full-arm execution makes every feasible
action observable; it does not by itself validate conditional confidence
bounds.}
\section{Retained decision theory}
\label{sec:b-theory}

This section records the decision-theoretic results the first study builds on.
They are \emph{retained, not newly claimed}: the ordinal-sufficiency theorem is
the known advantage-threshold result of the budgeted-triage literature in
ordinal form \citep{okati2021,fernandezloria2022,kamran2024}, and the
pessimistic gate is a standard one-step safe-improvement argument.  What the
paper tests is the behaviour of \emph{fitted} policies under matched
information, budget and feasibility --- a comparison the population-level
results below do not settle --- together with the operational question of how
action feasibility, verification cost and clean-run harm interact in
allocation.

One reading organises everything that follows.  The design is a horizon-one
decision process with a context-dependent action set: at a sampled checkpoint
the controller selects one intervention option from the feasible set, the frozen
continuation $\rho$ then runs to the horizon, and validated completion is the
terminal reward.  Every quantity the branch design estimates is therefore a
Bellman object of the baseline controller that never intervenes.  Writing
$V^\rho$ and $Q^\rho$ for its state- and action-value functions at the sampled
checkpoint, evaluated at the augmented state $X$ of \cref{sec:b-problem}, the
conditional regime value $\mu_g$ is an aggregate of $Q^\rho(\cdot,g)$, baseline
failure \emph{risk} is $1-V^\rho$, and verification advantage and actionable
value are advantages $A^\rho=Q^\rho-V^\rho$; a branch is a Monte-Carlo backup of
$Q^\rho$ from a restored state, and the availability mask is membership in an
option's initiation set.  \Cref{tab:bellman-dictionary} in \cref{app:b-notation}
gives the full correspondence.  The reading is used for two purposes only.  It
states the retained policy-ranking result in familiar terms --- \cref{thm:ordinal} is the exact
condition under which the failure-risk function $1-V^\rho$ ranks the advantage
function $A^\rho(\cdot,g_v)$; ranking by $V^\rho$ itself is the same comparison
with the sign reversed, and it is failure risk that a monitor reports --- and,
in \cref{sec:bellman}, it makes the one-step restriction a quantified choice
rather than a scoping sentence.  It is not a method: the learner of
\cref{sec:rl} remains a one-step full-information procedure that optimises a
masked policy score rather than regressing an action value, and no sequential
value iteration is proposed.  Nor is it an identification argument.  A Bellman
recursion is an accounting identity for whichever kernel generated the data; it
is the branch-data contract and the runtime qualification boundary
(\cref{sec:b-data}), not the recursion, that decide whether that kernel is the
one the branches executed.

What the retained results determine.  Each fixes a component of the design or a
prediction in \cref{sec:policy,sec:b-experiments}; none requires data.  The
first supplies the sharp null of the population-level comparison: exactly when
a baseline failure-risk score suffices to allocate a verification budget --- in
the reading above, exactly when $1-V^\rho$ ranks $A^\rho(\cdot,g_v)$
(\cref{thm:ordinal}).  The remaining results are the machinery that makes acting
on that answer legitimate --- a bound on the loss from registering a
representation (\cref{thm:refinement}), a pessimistic gated rule that cannot
make a running system worse (\cref{thm:safe,cor:proposal}), and a simultaneous
certification whose scope the first study narrows rather than inherits
(\cref{prop:certification}).

The results are individually elementary, and we state that plainly: the
contribution of this paper is not these theorems but the decision problem they
are applied to, together with the observation that
\emph{the one-step restriction makes their hypotheses verifiable rather than
assumed}.  A sequential safe-improvement guarantee of the kind in
\cref{thm:safe} would require uniform confidence validity over the states
reachable under the new policy, which is unknown and demands a positivity
assumption over a sequential action space.  Executing every feasible regime at a
single randomly selected checkpoint removes that problem for the
\emph{action}: conditional on a checkpoint being sampled, every regime feasible
there is executed, so action-assignment positivity holds by construction rather
than by assumption.  It does not promise that every cell of a partition is
sampled: a cell the registered sampling law gives zero or negligible mass
receives no branches, and the gate treats it as uncovered and uncertifiable
rather than as certified by default.

Two things must be kept apart.  Exhaustive execution supplies
\emph{coverage}: every feasible regime is executed at every eligible checkpoint,
so for the execution indicator $E_{ig}$ $\Prb(E_{ig}=1\mid A_{ig}=1)=1$.  It does
not supply \emph{statistical} validity: the uniform confidence bounds
\cref{thm:safe} requires are the gate family of \cref{eq:gate-family}, which
depends on a registered finite partition of the representation space, the sample
size, and a simultaneous union bound over the gate and certification families on
held-out calibration data.  \textbf{Full-arm execution makes every feasible
action observable; it does not by itself validate conditional confidence
bounds.}

Throughout, $\mu^{(j)}_g(z,m)=\E[Y^g_i\mid Z_i=z,M_i=m,J_i=j]$ denotes the
stratum-specific conditional regime value, defined where $g$ is available.


Throughout, $\mu^{(j)}_g(z,m)=\E[Y^g_i\mid Z_i=z,M_i=m,J_i=j]$ denotes the
stratum-specific conditional regime value, defined where $g$ is available.

\paragraph{Standing assumptions.}
The following are used throughout and are collected here so that each result can
name the ones it needs.  \Cref{as:baseline,as:unit} are enforced by the design;
\cref{as:frozen} by the protocol; \cref{as:continuity} is a regularity condition
whose failure mode is described below.

\begin{assumption}[Baseline availability]
\label{as:baseline}
$A_{i,g_0}=1$ at every eligible checkpoint: no intervention is always feasible.
This is what makes $G_j(d)$ in \eqref{eq:stratum-values} well defined, and what
makes \eqref{eq:actionable} a maximum over a non-empty set.
\end{assumption}

\begin{assumption}[One checkpoint per root; consistency under a regime]
\label{as:unit}
Exactly one checkpoint is sampled per root trajectory, and roots are independent
\emph{across clusters}, a cluster being the task template a root is instantiated
from.  For a policy $d$, the composite potential outcome satisfies
$Y^{d(Z_i,M_i)}_i=\sum_{g}\mathbf 1\{d(Z_i,M_i)=g\}\,Y^g_i$, so that
$\E[Y^{d(Z,M)}\mid Z,M,J]=\mu^{(J)}_{d(Z,M)}(Z,M)$.
\end{assumption}

The clustered form is the operative one, because the design instantiates several
roots from each template; independence at root resolution is not available and is
not assumed.  The two halves of \cref{as:unit} are used in different places, and
it is worth separating them.  Consistency is what
\cref{thm:refinement,thm:safe,cor:proposal} consume: those are conditional-mean
identities and tower-property arguments, and they hold whatever the dependence
between units.  Independence enters only through the confidence bounds priced in
\cref{prop:certification}, so every finite-sample count there is a count of
independent \emph{clusters} and reading it as a root count is anticonservative.

\begin{assumption}[Frozen comparator]
\label{as:frozen}
$d^\star\in\Dcal_{\mathrm{fixed}}$ and the candidate class $\Dcal$ are fixed
before the data used to construct confidence bounds are observed, so that the
advantage \eqref{eq:pointwise-advantage} is a fixed population quantity and the
coverage statements of \cref{prop:certification}, stated in full below, refer
to a fixed estimand.  \Cref{sec:policy} is authoritative about class membership,
derived gated outputs, and the split at which each object is frozen; reusing
calibration data to adapt a member would instead make the target data-dependent.
\end{assumption}

\begin{assumption}[Continuity of the ranking scores]
\label{as:continuity}
On the common-availability domain, $a_v(W)$ and the score used to allocate
capacity have continuous distributions, where $W=(Z,M)$.
\end{assumption}

\Cref{as:continuity} is needed only for \cref{thm:ordinal}, and only to make
capacity-$q$ selection unambiguous.  Under atoms, the upper level set need not
have probability exactly $q$; the registered analysis then uses randomised
tie-breaking at the boundary, both directions of \cref{thm:ordinal} continue to
hold with ``the same ordering'' read up to ties, and the equality case
identifies the selected sets only up to the atom.  We record this rather than
assume it away because a discrete risk score --- a bucketed monitor, say --- is
the likely case in practice.

\subsection{When does baseline failure risk rank verification value?}

The retained policy-ranking result concerns the cheapest identified comparison in the design:
verification against no intervention, on the domain where both are available.
Write $W=(Z,M)$ for the representation--mask pair of an eligible checkpoint and
let $a_v(w)=\mu_{g_v}(w)-\mu_{g_0}(w)$ on that domain.  Let $r$ be any baseline
failure-risk score on the same domain --- the registered risk score is a
function of $z$ alone, which is the special case $r(w)=r(z)$ --- and for a
capacity $q\in(0,1)$ let $S_q(s)$ denote the upper $q$-fraction of the eligible
population under score $s$.  The budget-matched risk-ranking regret is
\begin{equation}
 \mathcal R_{\mathrm{rank}}(q)
 =\E\big[a_v(W)\mid W\in S_q(a_v)\big]-\E\big[a_v(W)\mid W\in S_q(r)\big].
 \label{eq:rank-regret}
\end{equation}
Carrying the mask is not decoration.  Two checkpoints with identical $z$ but
different availability patterns are different conditioning cells and need not
share a verification advantage, so a regret written on $z$ alone would average
over a distinction the instrument was built to record.  The next statement is
\cref{thm:ordinal} with the mask carried explicitly.

\begin{theorem}[Ordinal sufficiency]
\label{thm:ordinal}
Under \cref{as:continuity}, $\mathcal R_{\mathrm{rank}}(q)\ge0$ for every $q\in(0,1)$, and
\begin{equation}
 \begin{aligned}
 \mathcal R_{\mathrm{rank}}(q)=0\ \text{ for every }q\in(0,1)
 &\quad\Longleftrightarrow\quad\\[-2pt]
 &r\ \text{and}\ a_v\ \text{induce the same ordering almost everywhere.}
 \end{aligned}
\end{equation}
\end{theorem}

\Cref{thm:ordinal} says that baseline failure risk $1-V^\rho$ ranks the
advantage function $A^\rho(\cdot,g_v)$ at every capacity if and only if the two
induce the same ordering almost everywhere; it converts a vague methodological
complaint into a hypothesis with a sharp null.  The sign matters and is not
cosmetic: the top-$q$ set under $r=1-V^\rho$ is the \emph{bottom}-$q$ set under
$V^\rho$, so ``$V^\rho$ ranks $A^\rho$'' is a different --- and, except in
degenerate cases, incompatible --- hypothesis.  The theorem depends on the score
only through the ordering it induces, so any strictly \emph{increasing}
transform of $r$ may be substituted freely; a strictly \emph{decreasing} one,
such as $V^\rho$ itself, may not.  The registered allocator spends its budget on
the highest failure risk, and that is the score the null is written on.  Write
\begin{equation}
 H_0:\quad\text{baseline failure risk is ordinally sufficient for verification advantage.}
 \label{eq:null}
\end{equation}
Under $H_0$, ranking checkpoints by predicted failure risk allocates a fixed
verification budget exactly as well as ranking them by the causal effect of
verifying, at every capacity, and the actionable-value framing of
\cref{def:actionable} buys nothing.  Rejecting $H_0$ is the precise sense in
which prediction and intervention value come apart.  This comparison requires
only $g_0$ and $g_v$, so it remains identified even if recovery semantics prove
incoherent.

\Cref{thm:ordinal} quantifies over an arbitrary score $r$; the registered test
does not.  It instantiates $r$ as one registered score: the cross-fitted
probability of baseline failure, fitted out of fold on the registered $Z$ alone,
on the same selection repetitions and template-grouped folds as the effect score,
from a learner library and hyperparameter grid frozen before any outcome
(\cref{app:e3-estimator}).  What is frozen before execution is the library and
the fitting procedure, not fitted weights.  A rejection therefore establishes
that \emph{that} score is ordinally insufficient on this population, not that no
risk score could be sufficient; because the two scores share data, folds and the
gradient-boosted family and differ in the label and in whether the mask enters, a
weak hand-built monitor is not an available explanation.  The paper claims only the former, and the scope of the
insufficiency framing should be read at that width.

The capacity-constrained shape has an adjacent empirical analogue: forecasting
improvements in storage placement produced monotone but sub-linear downstream
utility once the placement frontier was largely correct \citep{fleda2026}.
That observation motivates a descriptive transfer diagnostic; it is not
evidence for this theorem.

The null is not vacuous and the alternative is attainable.  Let $Z\sim
\mathrm{Unif}[0,1]$ with $r(z)=z$.  If $a_v(z)=z-\tfrac12$ the two orderings
agree and $\mathcal R_{\mathrm{rank}}\equiv0$.  If instead
$a_v(z)=\tfrac12-z$, then $S_q(r)=[1-q,1]$ while $S_q(a_v)=[0,q]$, and
\begin{equation}
 \mathcal R_{\mathrm{rank}}(q)
 =\Big(\tfrac12-\tfrac q2\Big)-\Big(\tfrac q2-\tfrac12\Big)=1-q>0
 \quad\text{for every }q\in(0,1).
 \label{eq:construction}
\end{equation}
Two populations with identical risk-score distributions therefore differ
maximally in whether that score is usable for allocation.  Which case obtains on
real checkpoints is an empirical question, and it is the one
\cref{sec:b-experiments} is organised around.

\paragraph{What the theory requires of the instrument.}
\Cref{eq:construction} also constrains the fault mechanisms of
\cref{sec:b-data}: under corruption detectable exactly when it is damaging
--- a substituted dependency, an altered manifest --- baseline risk and
verification advantage co-order, $H_0$ is not rejected, and the design has
measured its own injector.  Reaching the alternative requires a semantic fault
family --- an invalid inference drawn from a valid computation, with no trace
anomaly and no privileged rollback boundary.  The theory determines a
requirement on the instrument rather than inheriting one from it.

\subsection{Representations, representation loss, and why the mask is a policy input}

Write $\Dcal_H$ for the set of policies measurable with respect to $H_i$ and
$\Dcal_Z$ for those measurable with respect to $(Z_i,M_i)$; every policy in
either class selects a regime that is available at the checkpoint.  Let
$\mu^{(j)}_g(H)=\E[Y^g\mid H,J=j]$ and define the best achievable stratum values
\begin{equation}
 V^{*H}_j=\sup_{d\in\Dcal_H}V_j(d),
 \qquad
 V^{*Z}_j=\sup_{d\in\Dcal_Z}V_j(d),
 \qquad
 \mathrm{Ref}_j=V^{*H}_j-V^{*Z}_j .
 \label{eq:refinement}
\end{equation}

The representation question needs one definition.  Say that $Z$ is
conditionally-mean sufficient for the registered decision problem when
\begin{equation}
 \E[Y^g_i\mid H_i,J_i=j]=\E[Y^g_i\mid Z_i,M_i,J_i=j]
 \quad\text{for every feasible }g,
 \label{eq:conditional-mean-sufficiency}
\end{equation}
which is what would be required to claim no decision-relevant conditional-mean
loss relative to policies using the complete history.

\begin{theorem}[Representation loss]
\label{thm:refinement}
Assume \cref{as:unit}.  The feasible set $\Gset(H_i)$ is $M_i$-measurable.  Consequently
\begin{equation}
 V^{*H}_j=\E\Big[\max_{g\in\Gset(H)}\mu^{(j)}_g(H)\;\Big|\;J=j\Big],
 \qquad
 V^{*Z}_j=\E\Big[\max_{g\in\Gset(H)}\mu^{(j)}_g(Z,M)\;\Big|\;J=j\Big],
\end{equation}
and $\mathrm{Ref}_j\ge 0$.  If conditional-mean sufficiency
\eqref{eq:conditional-mean-sufficiency} holds, then $\mathrm{Ref}_j=0$.
Finally, for any representation $Z'$ with $\sigma(Z)\subseteq\sigma(Z')\subseteq
\sigma(H)$,
\begin{equation}
 \mathrm{Ref}_j\;\ge\;V^{*Z'}_j-V^{*Z}_j\;\ge\;0 .
 \label{eq:refinement-bound}
\end{equation}
\end{theorem}

The measurability clause is not a technicality.  By construction
$M_i=(A_{ig}:g\in\Gset)$ with $A_{ig}=\mathbf 1\{g\in\Gset(H_i)\}$, so
$\Gset(H_i)$ is recoverable from $M_i$ alone; without the mask, the maximisation
defining $V^{*Z}_j$ would range over a set the policy cannot identify.
\Cref{thm:refinement} is therefore the formal reason the mask is a policy input
rather than a nuisance.

\Cref{eq:refinement-bound} makes representation adequacy an empirical question.
$\mathrm{Ref}_j$ is not identified, because $\mu^{(j)}_g(H)$ is not estimable at
checkpoint resolution, but $V^{*Z'}_j-V^{*Z}_j$ is estimable from the same
branch data for any richer registered $Z'$: a representation audit returns a
lower bound on representation loss rather than a verdict.  This is the estimand
of the representation experiment in \cref{sec:b-experiments}.

\subsection{Actionable value}

\begin{definition}[Actionable value]
\label{def:actionable}
The actionable value of a checkpoint at stratum $j$ is
\begin{equation}
 \mathrm{AV}^{(j)}(z,m)=\max_{g\,:\,m_g=1}\big\{\mu^{(j)}_g(z,m)-\mu^{(j)}_{g_0}(z,m)\big\},
 \label{eq:actionable}
\end{equation}
the largest improvement any available regime achieves over no intervention.
\end{definition}

By \cref{as:baseline}, $g_0$ is an admissible element of the maximisation, so
$\mathrm{AV}^{(j)}\ge0$ pointwise, with equality exactly when no available regime improves
on continuing.  Actionable value therefore cannot detect harm on its own, and
should not be asked to: harm is a property of an \emph{individual} regime, not
of the best one, and the decision-relevant quantity is whether some available
$g$ carries a negative contrast at a checkpoint.  The
two are reported together, and neither substitutes for the other.

This is a substantive difference from formulations that measure advantage
against a baseline policy which may itself intervene; there the analogous
quantity can go negative, and a single number conflates ``nothing helps'' with
``acting hurts''.  Separating them is what allows the no-fault stratum to carry
its own endpoint.

Actionable value is the quantity a controller needs and a failure-risk score
does not supply: baseline risk $1-\mu^{(j)}_{g_0}(z,m)$ measures what happens if
nothing is done, $\mathrm{AV}^{(j)}$ measures what any available action can still change,
and \cref{thm:ordinal} establishes the exact condition under which the first
determines the second.

\subsection{Certification under simultaneous constraints}

A policy is eligible only if it certifies no-fault noninferiority
\eqref{eq:noninferiority}, the invalid-output bound
\eqref{eq:invalid-constraint}, and each of the $K$ cost components
\eqref{eq:cost-constraint}.  Write $\widehat{\mathcal C}$ for the resulting
certified set and $\mathcal C$ for the policies that truly satisfy all
constraints.  $\widehat{\mathcal C}$ is a random set whose selection event is not
controlled by per-constraint allocation without simultaneity over $\Dcal$.  The
$K+3$ bounds per member --- value, noninferiority, invalidity, and $K$ cost ---
form the \emph{certification family}, union-bounded on the same held-out data
with the gate family of \cref{eq:gate-family} defined in the next subsection, so
the price below covers both.  Derived deployment rules such as $d^\dagger$ and
$\mathrm{rl}_{\mathrm{constrained,certified}}$ of \cref{sec:policy} are not
members: their safety is the content of \cref{thm:safe,cor:proposal}, and their
value is reported from the untouched final split.

\begin{proposition}[Joint certification and its price]
\label{prop:certification}
Let $\Dcal$ be finite.  The protocol constructs one simultaneous family on the
held-out calibration data whose members are the certification bounds --- for
each member of $\Dcal$, one value, one no-fault noninferiority, one
invalid-output, and $K$ cost bounds --- together with the gate family of
\cref{eq:gate-family}, of combined size
\begin{equation}
 N=|\Dcal|(K+3)+F_{\mathrm{gate}} .
 \label{eq:combined-family}
\end{equation}
Suppose this combined family is simultaneously valid at level $1-\delta$.  Then
$\Prb(\widehat{\mathcal C}\subseteq\mathcal C)\ge 1-\delta$, and any policy
selected from $\widehat{\mathcal C}$ satisfies every constraint with probability
at least $1-\delta$ jointly; simultaneously, the gate-family members deliver the
uniform-validity hypothesis of \cref{thm:safe}.  Allocating $\delta/N$ to each
member suffices.  For a member $k$ whose contribution per independent unit is
bounded in an interval of length $R_k$ and which is averaged over $n_k$
independent units, a two-sided Hoeffding half-width at that allocation is
\begin{equation}
 \varepsilon_k(n_k)=R_k\sqrt{\frac{\log\!\big(2N/\delta\big)}{2n_k}},
 \label{eq:halfwidth}
\end{equation}
and this Hoeffding half-width falls below a true slack $s_k>0$ only when
\begin{equation}
 n_k\;\ge\;\frac{R_k^{2}\,\log\!\big(2N/\delta\big)}{2s_k^{2}} .
 \label{eq:sample-ceiling}
\end{equation}
When the independent units are clusters of unequal size with declared weights
$w_c$ summing to one, the same argument applied to the weighted mean of cluster
contributions gives $\varepsilon_k=R_k\sqrt{\log(2N/\delta)\sum_c w_c^{2}/2}$,
which reduces to \eqref{eq:halfwidth} with $n_k=1/\sum_c w_c^2$ at equal
weights.
\end{proposition}

Four things about \eqref{eq:halfwidth} are load-bearing and were previously
elided.  The range $R_k$ is member-specific: a completion rate in $[0,1]$ has
$R_k=1$, a paired contrast or a no-fault difference in $[-1,1]$ has $R_k=2$ and
needs \emph{four times} as many independent units for the same width, and a
cost component has whatever bounded range its registry ceiling fixes.  The
count $n_k$ is member-specific too: a gate bound on one cell, mask and stratum
is averaged over the units sampled into that cell, not over all roots, so
sparse cells have wide bounds and cells with no sampled units have none.  The
independent unit is the task template; the weighted form above is the one the
scaffold implements, and reading $n_k$ as a root count is anticonservative
whenever roots share a template.  Finally, \eqref{eq:halfwidth} is a planning
width for this particular bound: \eqref{eq:sample-ceiling} is neither a sample
size that every valid procedure needs nor a power guarantee.  A bound certifies
with high probability only when the slack exceeds roughly twice the half-width,
because the estimate itself fluctuates on that scale.  It is a worst-case bound, so a variance-adaptive
construction using pilot variance will be tighter; the allocation $\delta/N$ is
unaffected, since it concerns which bounds hold together, not how each is built.

\Cref{eq:sample-ceiling} is a planning width for a member of the
certification family, not a sample-size mandate.  The historical manuscript
priced a family with $N=291$ members at the registered values $\delta=0.05$,
$K=8$, $|\Dcal|=9$, $B=8$, $|\mathcal M|\le4$, $|\Gset|=3$; that family and its
count are a \emph{historical registration} and are not inherited by the first
study here.  The planned route is a small deployable policy family: fit the
advantage learner on training data, construct a small threshold/allocation family
on validation data, evaluate paired policy gains and no-injected-fault
degradation on independent calibration templates, and spend one explicitly
allocated simultaneous error budget over that family's endpoints.  A
bound-derived count is reported as a planning width for whichever member is
actually certified, with member-specific ranges and the independent unit being
the template rather than the root.  A design that cannot clear its planning
width coarsens the partition, reduces the policy family, or abandons the
certified claim in favour of the descriptive contract before, not after, the
confirmatory data are seen.

\subsection{Availability-masked pessimistic safe improvement}

Let $d^\star\in\Dcal_{\mathrm{fixed}}$ be the non-personalized comparator, frozen
before the certification data are observed.  Define the pointwise advantage of
an available regime relative to that comparator,
\begin{equation}
 \Delta^{(j)}_g(z,m)=\mu^{(j)}_g(z,m)-\mu^{(j)}_{d^\star(z,m)}(z,m),
 \label{eq:pointwise-advantage}
\end{equation}
which is defined wherever $A_g=1$, without a common-domain restriction, because
$d^\star(z,m)$ is available by construction; with cell-measurable policies the
cell-level value \eqref{eq:cell-advantage} agrees with it on each cell.  Let
$\mathrm{LCB}_\delta[\cdot]$ denote lower confidence bounds computed from the
held-out branch data.  The learned policy is defined as a \emph{gated deviation}
from the comparator, evaluated at the cell containing the checkpoint,
\begin{equation}
 d^\dagger(c)=
 \begin{cases}
  \displaystyle\argmax_{g\,:\,A_g=1}\;\min_{j\in\{0,1\}}\mathrm{LCB}_\delta\big[\Delta^{(j)}_g(c)\big],
   & \text{if }\displaystyle\max_{g\,:\,A_g=1}\;\min_{j\in\{0,1\}}\mathrm{LCB}_\delta\big[\Delta^{(j)}_g(c)\big]>0,\\[2ex]
  d^\star(c), & \text{otherwise.}
 \end{cases}
 \label{eq:gated-rule}
\end{equation}
where $d^\dagger(z,m)=d^\dagger(c(z,m))$ recovers the policy on the original
space.

\paragraph{The gate family.}
\Cref{thm:safe} below requires lower confidence bounds simultaneously valid over
the eligible support.  As written that is a function-class statement, so the
protocol registers, before calibration, a finite partition of the representation
space into $B$ cells, with $c(z,m)$ denoting the cell containing $(z,m)$.  The
policy class is restricted to cell-measurable policies, $d(z,m)=d(c(z,m))$, and
every quantity the gate uses is defined at cell resolution: the conditional
advantage given cell $c$ is
\begin{equation}
 \Delta^{(j)}_g(c)=\mu^{(j)}_g(c)-\mu^{(j)}_{d^\star(c)}(c),
 \qquad
 \mu^{(j)}_g(c)=\E\big[Y^g_i\mid c(Z_i,M_i)=c,\ J_i=j\big],
 \label{eq:cell-advantage}
\end{equation}
defined wherever $g$ is available.  This is not a patch: binning $Z$ into cells
is a coarsening of the registered representation, and the loss it induces is
exactly what \cref{thm:refinement} bounds, so the partition is a representation
choice the theory already prices.  The gate then evaluates one lower confidence
bound per cell, mask pattern, available regime, and stratum, so the family has at
most
\begin{equation}
 F_{\mathrm{gate}}\;\le\;B\,|\mathcal M|\,|\Gset|\,2 ,
 \label{eq:gate-family}
\end{equation}
members, where $\mathcal M$ is the set of mask patterns ($|\mathcal M|\le4$ by
\cref{as:baseline}) and $\Gset$ the registered regime set.  A union bound over
this family on the held-out calibration data delivers exactly the cell-level
uniform-validity hypothesis that \cref{thm:safe} consumes.  This is the gate
family.  It is distinct from, and union-bounded together with, the
certification family of \cref{prop:certification}; \cref{prop:certification} is
cited only for the latter, never for the gate bounds.

\begin{theorem}[Safe improvement, uniform in prevalence]
\label{thm:safe}
Assume \cref{as:baseline,as:unit,as:frozen}, and that with probability at least
$1-\delta$ the gate bounds are
simultaneously valid, $\mathrm{LCB}_\delta[\Delta^{(j)}_g(c)]\le
\Delta^{(j)}_g(c)$ for every available $g$, every cell $c$ in the registered
partition, and $j\in\{0,1\}$.  Then with probability at least $1-\delta$,
\begin{equation}
 V_j(d^\dagger)\ge V_j(d^\star)\ \text{ for }j=0,1,
 \qquad\text{and hence}\qquad
 V(d^\dagger;p)\ge V(d^\star;p)\ \text{ for every }p\in[0,1]
 \label{eq:safe-improvement}
\end{equation}
simultaneously.
\end{theorem}

Three features distinguish \eqref{eq:safe-improvement} from its sequential
antecedents.  The one-step performance difference is exact and consists of a
single term, so no occupancy measure or remaining-horizon factor appears.  The
uniform-validity hypothesis is over the registered partition, on which
full-arm execution supplies feedback for every feasible regime at every
\emph{sampled} checkpoint --- cells the sampling law does not reach stay
uncovered and are never deviated on --- and it is discharged by a union bound
over the gate family of \cref{eq:gate-family} ---
whereas a sequential design would need the hypothesis over a policy-dependent
reachable set, requiring a positivity assumption the design cannot enforce.  The
mask enters: the comparator is context-dependent, so \eqref{eq:pointwise-advantage}
is well defined without a common-domain restriction.

Two degenerate cases are informative rather than pathological.  If the deviation
set is empty the rule returns $d^\dagger=d^\star$ and
\eqref{eq:safe-improvement} holds with equality: the guarantee holds vacuously,
which is why the fallback rate is a reported endpoint.  If the uniform-validity
hypothesis fails the theorem says nothing --- it does not degrade gracefully ---
so bound-family coverage is checked directly on held-out branches and a
controller whose bounds fail is reported as unguaranteed rather than shipped
with a caveat.

Two design choices in \eqref{eq:gated-rule} are load-bearing.  Defining the
advantage against $d^\star$ rather than $g_0$ makes the fallback branch
contribute exactly zero rather than an unsigned quantity.  Taking the minimum
over $j$ before gating buys uniformity in $p$: the policy deviates only where
the action helps whether or not the registered fault was injected, so the
no-fault constraint \eqref{eq:noninferiority} is enforced pointwise.  A weaker
mixture gate is available when that is too conservative, at the cost of
uniformity.

\begin{corollary}[Guarantee interval]
\label{cor:interval}
Let $\Gamma_j=V_j(d^\dagger)-V_j(d^\star)$ and
$\Gamma(p)=p\Gamma_1+(1-p)\Gamma_0$.  Then $\{p\in[0,1]:\Gamma(p)\ge0\}$ is a
closed subinterval of $[0,1]$.  If $\Gamma_0$ and $\Gamma_1$ have strictly
opposite signs its finite endpoint is $-\Gamma_0/(\Gamma_1-\Gamma_0)$, the
break-even prevalence of \cref{app:estimands}.  Under the gate of
\eqref{eq:gated-rule} the interval is all of $[0,1]$.
\end{corollary}

\Cref{cor:interval} gives the break-even algebra of \cref{app:estimands} an
operational reading: it is not a descriptive calculation but the boundary of the
region over which the certificate holds.

The proof of \cref{thm:safe} never uses \emph{how} the candidate regime was
produced, so the guarantee extends to any proposal policy.

\begin{corollary}[Proposal-agnostic safe improvement]
\label{cor:proposal}
Let $d_{\mathrm{prop}}$ be a proposal policy that is (a) frozen independently of
the calibration data, (b) restricted to the registered feasible set,
$d_{\mathrm{prop}}(z,m)\in\Gset(z,m)$ at every eligible checkpoint, and (c)
depends only on pre-treatment information, in particular on availability
determined before treatment, and is cell-measurable.  Define the gated deviation
toward $d_{\mathrm{prop}}$,
\begin{equation}
 \begin{aligned}
 d^{\dagger}_{\mathrm{prop}}(c)&=
 \begin{cases}
  d_{\mathrm{prop}}(c), & \text{if }\displaystyle\min_{j\in\{0,1\}}\mathrm{LCB}_\delta\big[\Delta^{(j)}_{d_{\mathrm{prop}}(c)}(c)\big]>0,\\[1ex]
  d^\star(c), & \text{otherwise,}
 \end{cases}\\
 d^{\dagger}_{\mathrm{prop}}(z,m)&=d^{\dagger}_{\mathrm{prop}}(c(z,m)).
 \end{aligned}
 \label{eq:proposal-rule}
\end{equation}
Under the hypotheses of \cref{thm:safe}, with probability at least $1-\delta$,
\begin{equation}
 V_j(d^{\dagger}_{\mathrm{prop}})\ge V_j(d^\star)\ \text{ for }j=0,1,
 \qquad\text{and hence}\qquad
 V(d^{\dagger}_{\mathrm{prop}};p)\ge V(d^\star;p)\ \text{ for every }p\in[0,1]
 \label{eq:proposal-safe}
\end{equation}
simultaneously, regardless of the quality of $d_{\mathrm{prop}}$.
\end{corollary}

\Cref{thm:safe} is the special case of \cref{cor:proposal} with
$d_{\mathrm{prop}}(z,m)=\argmax_{g\,:\,A_g=1}\min_{j\in\{0,1\}}
\mathrm{LCB}_\delta[\Delta^{(j)}_g(c(z,m))]$.  The conditions on
$d_{\mathrm{prop}}$ are each load-bearing, and in different ways.  Restriction
to the registered feasible set is needed for \cref{thm:safe} itself: a regime
never branched has no bound to gate on, so a proposal that emits a patch outside
$\Gset$ at deployment time carries no guarantee.  Independence from the
calibration data is needed not by the proof --- the gate family of
\cref{eq:gate-family} is valid however the proposal was produced --- but by the
certification argument: a proposal fit on the calibration data could not be a
member of the pre-registered class $\Dcal$, so \cref{prop:certification} would
not certify its constraints and the target of inference would be data-dependent
(\cref{as:frozen}).  Pre-treatment availability keeps the choice a policy rather
than a post-treatment selection.  In short, (a) is a condition for the hypothesis
to be \emph{attainable}, not for the argument to be \emph{valid}; the proof
consumes only the validity event and the mask.  This is a certification result,
not an RL result, and it is what makes the design useful to the
harness-optimization literature of \cref{sec:related}: \emph{the proposal
proposes; the gate certifies}.

\paragraph{Dependency.}  \Cref{thm:safe} is void if branches do not begin from
an identical restored state.  Its hypotheses are guaranteed by the design only
insofar as restoration fidelity and branch isolation are demonstrated, which is
the object of the instrument gate in \cref{sec:b-experiments} and not something this
paper assumes.

%% --- M2: inserted from notes/mdp_framing/checkpoint_mdp.tex (argument order preserved) ---
\subsection{The design as a one-step Bellman backup}
\label{sec:bellman}

The reading that opened this section is now used for its second purpose: it
fixes what a sequential extension would have to add, and prices it.  Index the
eligible checkpoints of a root by $t=1,\dots,T$ and let $X_t$ be the augmented
state of \cref{sec:b-problem} at checkpoint $t$ --- history, latent stratum,
remaining budget, remaining horizon.  The primitive action is the
\emph{intervention option} of \cref{tab:option-boundaries}, not the complete
regime: the option initiates where $A_{ig}=1$, runs its registered action and
frozen response rule, and terminates at its postcondition, after which $\rho$
resumes.  The complete regime is that option \emph{plus} $\rho$ to the horizon,
which is why a complete regime cannot also serve as the next-checkpoint
primitive.  Options have variable duration and consume budget, so the process
over checkpoints is properly semi-Markov; it reduces to a finite-horizon MDP
because the protocol fixes the eligible grid and because elapsed time and
remaining budget are carried inside $X_t$.  With that reduction the kernel is
\begin{equation}
 P_{t,g}(x'\mid x)=\Prb\{X_{t+1}=x'\mid X_t=x,\ \text{execute option }g,\ \text{then }\rho
 \text{ to the next decision point}\},
 \label{eq:checkpoint-kernel}
\end{equation}
for $g\in\Gset(x)$, with terminal reward $Y=R(X_{T+1})\in\{0,1\}$ and absorbing
states for early termination.  The baseline controller selects $g_0$ at every
checkpoint (\cref{as:baseline} makes it admissible).  Write $V^\rho_t$,
$Q^\rho_t$, and $A^\rho_t=Q^\rho_t-V^\rho_t$ for its value, action-value, and
advantage functions, so that
$Q^\rho_t(x,g)=\sum_{x'}P_{t,g}(x'\mid x)V^\rho_{t+1}(x')$ is the value of a
single deviation to $g$ at $(t,x)$ with future choices at the baseline.  A
branch restored at $(t,x)$ executes the option and then $\rho$, so its potential
outcome has mean $Q^\rho_t(x,g)$: with $\tau$ the sampled checkpoint,
\begin{equation}
 \begin{aligned}
 \mu^{(j)}_g(z,m)&=\E\big[Q^\rho_\tau(X,g)\mid Z=z,M=m,J=j\big],\\
 1-\mu^{(j)}_{g_0}&=1-\E[V^\rho_\tau(X)\mid\cdot],
 \qquad
 a_v=\E[A^\rho_\tau(X,g_v)\mid\cdot].
 \end{aligned}
 \label{eq:bellman-dictionary}
\end{equation}
Two aggregation points are load-bearing.  Conditioning on $J=j$ conditions on a
\emph{component of the state} rather than refining a history-measurable value:
had $Q^\rho$ been written on $H$ alone, the kernel would still depend on $j$ and
the value would have to be re-indexed as $Q^{\rho,j}$ before \eqref{eq:bellman-dictionary}
could be read.  And the conditioning law in \eqref{eq:bellman-dictionary} is the
registered root-and-checkpoint sampling law, so $\mu^{(j)}_g$ describes the
design's own population; a deployment mixture is formed from the stratum values
by \eqref{eq:prevalence-value} and by no other route.

Baseline failure risk is a \emph{state value}; verification advantage and
actionable value \eqref{eq:actionable} are \emph{advantages}; and
\cref{thm:ordinal} is the exact condition under which failure risk ranks the
advantage function.  \Cref{thm:refinement} is the value loss of the state
abstraction $Z=f(H)$ under one step of value iteration, with
\eqref{eq:conditional-mean-sufficiency} the condition that the abstraction be
$Q^\rho$-irrelevant; note that both $V^{*H}_j$ and $V^{*Z}_j$ are expectations
of maxima of \emph{conditional means}, and that
$\E[\max_g\mu^{(j)}_g(H)]\le\E[\max_g Q^\rho_\tau(X,g)]$ with the gap the
residual stochasticity a maximum of conditional means correctly declines to
claim.  \Cref{thm:safe} is pessimistic value iteration at horizon one on the
registered partition, and the gate family \eqref{eq:gate-family} is the size of
that tabular action-value table.  The branch design is depth-one sparse sampling
under generative-model access --- reset to a state, sample its successors --- at
the checkpoints the design samples and nowhere else, and E0 tests whether that
access exists rather than assuming it.  The learner of \cref{sec:rl} is a masked
policy-score learner: it maximises the estimated expected advantage
\eqref{eq:rl-objective} over a masked softmax and never regresses $Q^\rho$.  It
needs no bootstrap term for the same reason a fitted-$Q^\rho$ regression would
not: at horizon one the branch is the backup.  A fitted regression of the
outcome on the controller's observation $(z,m)$ and the regime is a coherent
\emph{alternative} deployable learner and is proposed as such in
\cref{app:b-analysis}; a regression on the full state $X$ would estimate
$Q^\rho(X,g)$ itself but is an oracle diagnostic, since $X$ carries the latent
stratum.  Neither is what \cref{sec:rl} implements, and neither is a member of
the registered class $\Dcal$.

What the reading buys is a formula for the sentence that opens this section.
Let $d$ be a cell-measurable rule with $d(c)\in\Gset(c)$, let $d^\rho_t$ be
the state distribution of $\rho$ at checkpoint $t$, and define
\begin{equation}
 \begin{aligned}
 \bar E_t(d)&=\E_{x\sim d^\rho_t}\big[A^\rho_t(x,d(x))\big],\\
 \alpha^{\mathrm{eff}}_t(d)&=\E_{x\sim d^\rho_t}
   \Big[\tfrac12\big\|P_{t,d(x)}(\cdot\mid x)-P_{t,g_0}(\cdot\mid x)\big\|_1\Big]
 \;\le\;\Prb_{x\sim d^\rho_t}\{d(x)\neq g_0\}.
 \end{aligned}
 \label{eq:transfer-terms}
\end{equation}
$\bar E_t$ is the mean advantage of $d$ over $g_0$ at checkpoint $t$, and
$\alpha^{\mathrm{eff}}_t$ is the kernel deviation the rule induces, bounded by
its intervention rate.  The bound is at $X$ resolution: a verification that
passes and continues still decrements the budget and appends the verifier's
output to $H_{t+1}$, so its kernel deviation is not zero.  Reducing
$\alpha^{\mathrm{eff}}$ for verification to the intervention rate times the
rate at which the response rule fires requires the additional assumption that
the pass branch changes only $V^\rho$-irrelevant coordinates of $H$ --- an
assumption of the same shape as \eqref{eq:conditional-mean-sufficiency}, to be
stated and audited rather than presumed.  Because $d$ is cell-measurable, the
window averages $\bar E=\frac1T\sum_t\bar E_t=\sum_c\Prb(c)\,\Delta_{d(c)}(c)$
and $\bar\alpha$ are identified by the design.

\begin{proposition}[Single random-checkpoint deployment]
\label{prop:l0}
Let $\mathcal W$ be the eligible window realised along a $\rho$-trajectory and,
conditional on $\mathcal W$, let $\tau$ be uniform on $\mathcal W$.  Deploy $d$
at $\tau$ and $g_0$ elsewhere.  Then
$V^{\mathrm{L0}}(d)-V^\rho=\E\big[\tfrac1{|\mathcal W|}\sum_{t\in\mathcal W}A^\rho_t(X_t,d(X_t))\big]
=\sum_c\Prb(c)\,\Delta_{d(c)}(c)$ exactly, where $\Prb(c)$ is the cell mass
under the design's sampling of $(\text{root},\tau)$, and under the gate of
\eqref{eq:gated-rule} this is non-negative.
\end{proposition}

\begin{proof}
Eligibility at $t$ is determined by $H_t$, and before $\tau$ the two policies
agree, so conditional on $\mathcal W$ and $\tau=t$ the state $X_t$ has $\rho$'s
law.  The value difference is then $A^\rho_t(X_t,d(X_t))$ in expectation; average
over $\tau$ given $\mathcal W$ and then over roots.  The cell form follows
because $d$ is constant on cells and the design samples $(\text{root},\tau)$
exactly this way.
\end{proof}

\Cref{prop:l0} is the trajectory-level content of \cref{thm:safe}: the
certificate covers a controller that intervenes once, at a checkpoint drawn as
the design draws it.  Online, that deployment is implementable as stated only
when the window is known in advance (a fixed checkpoint grid); a stopping-time
deployment adds a selection term the design does not estimate.  A controller
that acts at every checkpoint is a different object.

\begin{proposition}[Transfer to repeated deployment]
\label{prop:transfer}
Let $\pi_d$ select $d$ at every checkpoint, and let
$\mathrm{sp}_t(d)=\sup A^\rho_t(\cdot,d(\cdot))-\inf A^\rho_t(\cdot,d(\cdot))$
over the union of the supports of $d^{\pi_d}_t$ and $d^\rho_t$, both read at
$X$ resolution.  Then
\begin{equation}
 V^{\pi_d}-V^\rho
 =\sum_{t=1}^T\E_{x\sim d^{\pi_d}_t}\big[A^\rho_t(x,d(x))\big]
 \;\ge\;\sum_{t=1}^T\Big(\bar E_t(d)-\mathrm{sp}_t(d)\,\min\Big\{1,\sum_{s<t}\alpha^{\mathrm{eff}}_s(d)\Big\}\Big),
 \label{eq:transfer-bound}
\end{equation}
and since $\mathrm{sp}_t\le2$ and $\alpha^{\mathrm{eff}}_s\le\alpha_s$,
$V^{\pi_d}-V^\rho\ge T\,\bar E(d)-2T(T-1)\,\bar\alpha(d)$ in terms of window
averages alone.
\end{proposition}

\begin{proof}
The identity is the performance-difference lemma \citep{kakade2002}, by telescoping
$V^{\pi_d}_t-V^\rho_t$ along $\pi_d$'s occupancy.  For the bound, write
$d^{\pi_d}_{t+1}-d^\rho_{t+1}=(d^{\pi_d}_t-d^\rho_t)P_{\pi_d}+d^\rho_t(P_{\pi_d}-P_\rho)$;
the first term has $\ell_1$ norm at most $\|d^{\pi_d}_t-d^\rho_t\|_1$ and the
second at most $2\alpha^{\mathrm{eff}}_t$, so the total-variation distance
between the two occupancies at $t$ is at most
$\min\{1,\sum_{s<t}\alpha^{\mathrm{eff}}_s\}$ (the simulation lemma,
\citealp{kearnssingh2002}).  For any
$f$ and distributions $p,q$, $|\E_pf-\E_qf|\le\mathrm{span}(f)\,\mathrm{TV}(p,q)$
with the span over the union of supports; apply it to
$f=A^\rho_t(\cdot,d(\cdot))$ at each $t$.
\end{proof}

The per-checkpoint form of \eqref{eq:transfer-bound} is attained: in every
case of the two-state example below its right-hand side equals
$V^{\pi_d}-V^\rho$ exactly, so the constant $\mathrm{sp}_t$ cannot be
improved in general.  The window-averaged form is a relaxation and is not
attained; on random checkpoint MDPs it runs at about $-11$ against an actual
improvement of about $+0.1$, and it should be read only as the form
expressible in window averages.  The certificate transfers to repeated
deployment when
\begin{equation}
 \bar E(d)\;>\;\mathrm{sp}_{\max}(d)\,(T-1)\,\bar\alpha(d),\qquad \mathrm{sp}_{\max}(d)=\max_t\mathrm{sp}_t(d),
 \label{eq:transfer-condition}
\end{equation}
the mean certified one-step advantage exceeding the largest span of the
advantage times the deviation accumulated over the window.  One might object that the
worst-case span of $2$ is what makes this vacuous, and the objection is
correct as far as it goes: on random checkpoint MDPs the realised span of
$A^\rho$ is about $0.3$, the window-averaged form with the constant $2$ is
informative (above $-1$) in under $2\%$ of draws, and the span-refined form
in about $95\%$.  What the refinement does not do is make the condition
certifiable from the design.  The span is over the states the \emph{controller}
reaches, which include post-intervention states from which depth-one branching
never re-branches: the design observes $Y$ along those continuations but never
estimates $A^\rho$ there.  Replacing $\mathrm{sp}_t$ by the span over $\rho$'s
support --- the only span the branches identify --- yields a quantity that is
not a bound: in the example below with $\varepsilon=0$ it certifies an
improvement of $+0.05$ while the controller loses $0.05$.  Even with the
correct span the refined bound certifies improvement in well under $10\%$ of
random draws.  So the one-step restriction is not an artifact of a loose
constant.  It is the statement that one scalar --- the span of the one-step
advantage over post-intervention states --- is decision-relevant and
unmeasured, and that is the same coverage gap a sequential optimum would face:
$V^*$ needs backups at post-intervention states, and pessimistic value
iteration under a coverage assumption is the known remedy.  A depth-two
branching pilot would measure that scalar directly; it is a new instrument
requirement, not part of this submission.

\begin{remark}[Transfer fails at the registered scale]
\label{rem:transfer-fails}
Let two states $a,b$ share a cell.  Under $\rho$, $a$ leads to $b$ with
probability $\varepsilon=0.05$ and otherwise to a state $x$ with success
probability $0.85$; from $b$ the baseline succeeds with probability $0.90$.
Verification at $a$ routes to $b$; verification at $b$ succeeds with
probability $0.80$.  Then $A^\rho(a,g_v)=+0.0475$, $A^\rho(b,g_v)=-0.100$,
the cell's verification advantage is $+0.0405$, and the gated rule verifies on
the cell.  $V^\rho=0.8525$, $V^{\mathrm{L0}}=0.8738$, and the repeated
controller scores $0.8000$: every advantage and the loss are at the scale of
the no-fault margin $\eta=0.05$.  (With $x$ succeeding at $0.50$ and
verification at $b$ escalating, the same construction gives $V^\rho=0.52$,
$V^{\mathrm{L0}}=0.69$ and $V^{\pi_d}=0$.  Completion is binary, so absolute
harm can never exceed $1$ and is bounded here by $V^\rho=0.52$; what the variant
shows is that the loss is not confined to the $\eta$ scale but can consume the
entire baseline value.)  State $b$ carries $5\%$ of $\rho$'s mass and all of the
controller's.  That is a \emph{coverage} failure --- the density ratio between
the two occupancies is $20$ on the state that decides the outcome --- and not
yet the strict positivity failure named at the start of this section, because
$b$ is visited under $\rho$ and could in principle be branched from.  Now
set $\varepsilon=0$, so that $b$ is a post-intervention state $\rho$ never
visits, which \emph{is} a support failure: no amount of depth-one sampling
under $\rho$ reaches it.  The cell's advantage is $A^\rho(a,g_v)=+0.05$, the gate fires,
$V^{\mathrm{L0}}-V^\rho=+0.025$, and the repeated controller loses $0.05$.
The span of $A^\rho$ over $\rho$'s support at the second checkpoint is $0$
(only $x$ is visited, and $d(x)=g_0$), so the on-support version of
\eqref{eq:transfer-bound} evaluates to $+0.05$ and certifies an improvement
that does not exist, while the union-support form gives $-0.05$, which is
exact.  In random checkpoint MDPs the failure is rare and small: at the
registered seed the repeated cell-level rule is worse than $\rho$ in $0$ of
$1{,}500$ draws without post-intervention states and in $2$ of $1{,}500$ with
them (losses of $3\times10^{-5}$ and $3\times10^{-3}$, in both cases with
mass routed onto states $\rho$ never visits), and the on-support form of the
bound was never violated by chance.  The failure is structural, not typical,
and the bound is what separates the two.
\end{remark}
